\documentclass[12pt, oneside]{book}
\usepackage{lmodern} %fuentes
\usepackage[T1]{fontenc} %fuentes
\usepackage[spanish]{babel} %idoma
\addto\captionsspanish{\renewcommand{\listtablename}{Índice de tablas}}
\addto\captionsspanish{\renewcommand{\tablename}{Tabla}}
\usepackage{caption, subcaption} %figuras
\usepackage{graphicx}
\usepackage{array}
\usepackage{textcomp}  % 
\usepackage[width=155mm,top=25mm,bottom=25mm, headheight=15pt,]{geometry} %Tamaño de la Pagina
\usepackage{fancyhdr}
\usepackage{url}
\pagestyle{fancy} % Configura el estilo de las páginas
\fancyhf{} % Limpia las cabeceras y pies de página
\fancyfoot[R]{\thepage} % Número de página en el pie, alineado a la derecha
\usepackage[backend=biber,style=numeric,sorting=none]{biblatex}
\addbibresource{Referencias.bib}

\usepackage{pdflscape}
\usepackage{booktabs}
\usepackage{array}
\usepackage{caption}


%Abreviaturas
\usepackage[utf8]{inputenc}
\usepackage{array}
\usepackage{tabularx}
\newcolumntype{L}{>{\raggedright\arraybackslash}X}
\newcommand{\abbrentry}[3]{%
  \textbf{#1} – #2\nobreak\dotfill\nobreak#3%
}


% Configuración de líneas
\renewcommand{\headrulewidth}{0.4pt} % Línea en el encabezado
\renewcommand{\footrulewidth}{0.4pt} % Línea en el pie de página
\fancypagestyle{plain}{
    \fancyhf{}
    \fancyfoot[R]{\thepage} % Número de página en el pie, alineado a la derecha
    \renewcommand{\headrulewidth}{0.4pt} % Línea en el encabezado en el estilo plain
    \renewcommand{\footrulewidth}{0.4pt} % Línea en el pie de página en el estilo plain
}
\setlength{\parskip}{1em} % Ajusta el espacio entre párrafos
\setcounter{secnumdepth}{3}
\setcounter{tocdepth}{3}
% Paquete para manejar bibliografías
%\usepackage{natbib}
\usepackage{amsmath}
\usepackage{siunitx} % Opcional, para formatear unidades

\usepackage{soul} %Opcional, para hacer subrayado

\usepackage{tikz}
\usetikzlibrary{positioning, arrows}
\usepackage{colores-tesis} % paleta de las gráficas (carga pgfplots)

%
\usepackage[
    colorlinks=true,        % enlaces con color (no recuadros)
    linkcolor=blue,         % color para enlaces internos
    urlcolor=cyan,          % color para urls
    linktoc=all,            % hace clickeables tanto el texto como el número de página
    hidelinks               % opcional: quita colores y bordes, si prefieres “links invisibles”
]{hyperref}

\hypersetup{
    pdftitle={Título de tu Tesis},
    pdfauthor={Tu Nombre},
}

\begin{document}
	\begin{titlepage}
		\setlength{\parindent}{0pt} \setlength{\parskip}{0pt}
	
		\begin{center}
			\vfill 
			
			\begin{minipage}{\textwidth}
				
				\newcolumntype{V}{>{\centering\arraybackslash} m{.17\textwidth} }
				\newcolumntype{C}{>{\centering\arraybackslash} m{.56\textwidth} }
				
				\begin{center}
					\includegraphics[width=4.5cm]{UNAM_LOGO2.png}	
				\end{center} 	
				\begin{center}
					\Huge\textbf{Universidad Nacional Aut\'onoma de M\'exico}
				\end{center} 
				
			\end{minipage}
				
			\vfill
			
			\begin{minipage}{0.7\textwidth}
				\begin{center}
					\large Programa de Maestría en Ciencia e Ingeniería de la Computación
				\end{center}
			\end{minipage}
			% 
			\begin{center}
				
				\vfill
				%\medskip \rule{.9\textwidth}{2pt}
				\vfill
				\large Detección de anomalías en redes de monitoreo atmosférico mediante aprendizaje profundo
	%{\Large \bfseries  \title   }
				
				%\medskip \rule{.9\textwidth}{2pt}
			\end{center}
			
			\begin{center}
				\vspace*{1cm}
				{\huge T E S I S}\\
				\vspace*{1cm}
				que para optar por el grado de \\ \bigskip
				{\large {Maestro en Ciencia e Ingenier\'ia de la Computación}} \\ \bigskip
				
				%PRESENTA:\\ \medskip
				\begin{center}
					PRESENTA: \\
					\large José Emiliano Lili Beltrán
				\end{center}
				
				\bigskip
				
				TUTOR: \\
				Dr. Jos\'e Jaime Camacho Escoto, Facultad de Ingenier\'ia\\ UNAM
				%Coadvisors
				%\emph{Coasesores:}
				%\coadvisor1 \\ \coadvisor2 \\ \coadvisor3 
			\end{center}
			
			\vfill
			
			\begin{center}
				{Ciudad Universitaria, CD. MX. Junio, 2025}\\
			\end{center}
			\cleardoublepage
		\end{center}
	\end{titlepage}




% =====================
% Agradecimientos
% =====================
\chapter*{Agradecimientos}
\addcontentsline{toc}{chapter}{Agradecimientos}

AQUI VAN LOS AGRADECIMIENTOS

% =====================
% Índices
% =====================
\tableofcontents % Índice general
\listoffigures % Índice de figuras
\addcontentsline{toc}{chapter}{Índice de figuras}
\listoftables % Índice de tablas
\addcontentsline{toc}{chapter}{Índice de tablas}


\chapter*{Lista de abreviaturas}
\begin{tabularx}{\linewidth}{@{}Lr@{}}
  \abbrentry{NOM}{Norma Oficial Mexicana}{}\\
  \abbrentry{IoT}{Internet de las Cosas}{}\\
  \abbrentry{LoRaWAN}{Long Range Wide Area Network}{}\\
  \abbrentry{LoRa}{Long Range}{}\\
  \abbrentry{AES CTR}{Advanced Encryption Standard - Counter Mode}{}\\
  \abbrentry{ED}{End Device}{}\\
  \abbrentry{GW}{Gateway}{}\\
  \abbrentry{Plaintext}{Texto Plano}{}\\
  \abbrentry{Ciphertext}{Texto Cifrado}{}\\
  \abbrentry{MiTM}{Man-in-the-Middle}{}\\
  \abbrentry{MIC}{Message Integrity Code}{}\\
  \abbrentry{ML}{Machine Learning}{}\\
\end{tabularx}

% =====================
% Resumen y Abstract
% =====================
\chapter*{Resumen}
\addcontentsline{toc}{chapter}{Resumen}



\chapter*{Abstract}
\addcontentsline{toc}{chapter}{Abstract}


% =====================
% Capítulos
% =====================
% =====================
% Capítulo 1
% =====================
% =====================
% Capítulos (cada uno en su archivo)
% =====================
% =====================================================================
% Capítulo 1. Introducción                         Presupuesto: 7 páginas
% ---------------------------------------------------------------------
% Este capítulo es el «contrato» de la tesis: lo que prometas aquí (problema,
% hipótesis, objetivos) es exactamente lo que el capítulo 7 debe responder.
% Debe contestar:
%  - ¿Qué es lo que está en juego? (decisiones que dependen del dato)
%  - ¿Qué falla hoy? (valores alterados que siguen siendo plausibles)
%  - ¿Qué no han resuelto otros? (la brecha, con cita)
%  - ¿Qué propones y cómo sabrás si funcionó? (hipótesis medible)
% Texto del borrador de Overleaf conservado tal cual. Las marcas «% REVISAR»
% señalan lo que ya no coincide con lo que se hizo (dos etapas; NERAD).
% =====================================================================
\chapter{Introducción}

La adopción de las tecnologías de \textit{IoT} ha permitido la interconexión de dispositivos, el sensado, el monitoreo, la automatización y la transmisión de datos en entornos reales. Estas tecnologías empezaron a desempeñar un papel crítico en el funcionamiento de las industrias y las ciudades. Esto evidenció la necesidad de evolucionar hacia tecnologías de comunicación eficientes, de gran alcance y de bajo consumo energético.

Para atender esta necesidad, se empezaron a adoptar las redes \textit{LPWAN}. Estas redes encontraron aplicación en diversos sectores industriales, permitiendo el monitoreo de variables como la temperatura, la humedad y el pH, entre otras.

En 2009 se empezó a desarrollar una tecnología de modulación que permite la comunicación inalámbrica de bajo consumo y largo alcance, diseñada específicamente para entornos que requieren tecnologías de IoT en ciudades o industrias, llamada \textit{LoRa}. Sobre esta base se creó \textit{LoRaWAN}, un protocolo de red que permite transmitir pequeños volúmenes de datos con bajo consumo energético para su procesamiento por parte de servicios o aplicaciones, lo que apoya el monitoreo e incluso la toma de decisiones.

A lo largo de su evolución, \textit{LoRaWAN} ha incorporado distintos mecanismos orientados a proteger la comunicación, incluyendo el cifrado del \textit{FRMPayload} y mecanismos de integridad como el \textit{Message Integrity Code} (MIC). La versión 1.1 incorporó mejoras adicionales en materia de seguridad; sin embargo, la confidencialidad y la integridad del mensaje se proporcionan mediante mecanismos distintos. El cifrado del \textit{payload} se basa en operaciones XOR con un flujo de clave derivado de AES, por lo que una modificación realizada sobre un bit del texto cifrado puede reflejarse en la misma posición del texto claro una vez descifrado.

Esta propiedad permite considerar ataques de \textit{bit flipping} bajo un modelo de amenaza en el que un adversario tiene acceso al mensaje cifrado después de la verificación de integridad a nivel de red y antes de su procesamiento por la aplicación. Si además conoce la estructura utilizada para codificar el \textit{payload}, puede modificar posiciones específicas sin necesidad de conocer la clave de cifrado. Como consecuencia, una medición transmitida puede ser alterada y producir un valor distinto al original, cuya repercusión dependerá tanto de la posición modificada como del contexto en el que dicho dato sea utilizado.

A partir de este escenario, resulta relevante analizar no sólo la posibilidad de detectar una modificación, sino también el daño que ésta puede producir sobre la aplicación y las condiciones bajo las cuales una alteración puede permanecer difícil de distinguir del comportamiento legítimo de las mediciones.

% REVISAR: falta el puente hacia la etapa 2: la anomalía puede venir de un ataque,
% una falla o una descalibración, y el operador necesita saber QUÉ ESTACIÓN revisar.
\section{Definición del problema}

Con el aumento del despliegue de infraestructuras que utilizan \textit{LoRaWAN} como medio de comunicación, estas infraestructuras permiten transmitir una variedad de mediciones según la aplicación ya sea desde el nivel de humedad, la concentración de un gas o el pH de una ubicación. Estas mediciones pueden ser empleadas por aplicaciones que apoyan la toma de decisiones que pueden afectar procesos productivos o decisiones a escala urbana. 

En \textit{LoRaWAN}, la confidencialidad del \textit{FRM\textit{payload}} y la integridad del mensaje se proporcionan mediante mecanismos distintos. El cifrado del \textit{\textit{payload}} se basa en operaciones XOR con un flujo de clave derivado de AES, por lo que una modificación realizada sobre un bit del texto cifrado se refleja en la posición correspondiente del texto claro después del descifrado, aun cuando el atacante no conozca la clave. Por otra parte, la integridad de la trama se protege mediante el \textit{Message Integrity Code} (MIC). Sin embargo, trabajos previos han señalado que, en el segmento comprendido entre el \textit{Network Server} y el \textit{Application Server}, la comprobación de integridad no se preserva de extremo a extremo, lo que permite considerar un modelo de amenaza en el que un adversario con conocimiento de la estructura del \textit{\textit{payload}} puede modificar posiciones específicas del mensaje cifrado mediante un ataque de \textit{bit flipping}.

A pesar de que los ataques de \textit{bit flipping} han sido documentados y se han propuesto mecanismos para su detección, como el ICM presentado en \cite{alizadeh2023bitflipping}, su evaluación se orienta principalmente a distinguir entre mensajes legítimos y mensajes alterados. Para el caso de \textit{LoRaWAN}, los autores emplean el conjunto de datos \textit{Bike Sharing} de UCI y, adicionalmente, evalúan su propuesta sobre un conjunto de datos real de una red \textit{LoRaWAN}, recopilado mediante nueve \textit{gateways} en un entorno urbano. Las muestras atacadas son generadas mediante la modificación aleatoria de uno o más bits. Sin embargo, este procedimiento no caracteriza de manera sistemática el efecto asociado con la posición específica del bit modificado, ni relaciona dicha posición con el impacto producido sobre la magnitud representada y con su detectabilidad.

Esto conduce a plantear si es posible caracterizar el daño producido por los ataques de \textit{bit flipping} en función de la posición modificada y, al mismo tiempo, cuantificar su detectabilidad cuando los datos presentan variabilidad legítima. El problema no se limita únicamente a determinar si un mensaje fue alterado, sino a analizar si existen posiciones o condiciones en las que una modificación puede producir un impacto relevante sobre la decisión de la aplicación sin resultar fácilmente distinguible del comportamiento legítimo de las mediciones.

% REVISAR: la hipótesis sigue siendo válida. Falta hacerla medible: ¿mejorar respecto a
% qué (medición aislada, PARAAD-SA)? ¿con qué métrica y con qué tasa de falsas alarmas?
\section{Hipótesis}

Es posible detectar, con un alto grado de confiabilidad, anomalías en la información reportada en redes de sensores utilizando el historial reciente de cada estación y el de sus estaciones vecinas sin incurrir en un exceso de falsos positivos.

% MAPA OBJETIVO -> DÓNDE SE CUMPLE -> DÓNDE SE RESPONDE
% Hipótesis ........ prueba: OE3 (con/sin vecinas) y OE4 ......... cap. 7, sec. Respuesta a la hipótesis
% OE1 bit flipping ...... cap. 4 (método y resultados) ............ cap. 7
% OE2 diseño ............ cap. 5, sec. Arquitectura ................ cap. 7
% OE3 aporte vecinas .... cap. 6, sec. Elección del diseño (2022) .. cap. 7
% OE4 estado del arte ... cap. 6, secs. Una estación y Despliegue .. cap. 7
% OE5 dónde falla ....... cap. 6, secs. Coordinado y Sesgo ......... cap. 7
% Términos que se definen en el cap. 2: falsa alarma (= falso positivo en una
% estación-día), recall, estación vecina (cap. 5: las 4 más cercanas).
\section{Objetivos}

\subsection{Objetivo general}

Diseñar un detector de anomalías basado en modelos de aprendizaje profundo que combine la historia temporal de cada estación con la de sus estaciones vecinas, con alta capacidad de detección y una tasa de falsas alarmas acotada, en una red de sensores de ozono.

\subsection{Objetivos específicos}
\begin{itemize}
    \item[$\bullet$] Caracterizar el daño y la detectabilidad de los ataques de \textit{bit flipping} en LoRaWAN 1.1~\cite{loraalliance2017lorawan11} según la posición del bit modificado, a partir del cambio de banda de calidad del aire definido en la NOM-172-SEMARNAT-2023~\cite{nom172_semarnat_2024}.

    \item[$\bullet$] Diseñar un detector de anomalías que combine la historia reciente de la estación evaluada con la de sus estaciones vecinas e incorpore un filtro que excluya a las estaciones vecinas sospechosas de estar alteradas.

    \item[$\bullet$] Determinar cuánto aporta a la detección usar a las estaciones vecinas como referencia, frente a usar sólo la historia propia de la estación, con el mismo umbral de falsas alarmas.

    \item[$\bullet$] Evaluar el detector propuesto frente a un detector del estado del arte, mediante el \textit{recall} y los días de falsa alarma al año por estación, en un banco de pruebas y en un periodo de despliegue.

    \item[$\bullet$] Identificar las condiciones en las que el detector propuesto pierde capacidad de detección, como la alteración simultánea de estaciones vecinas o un sesgo sostenido en el tiempo.
\end{itemize}

% REVISAR: describe el esquema anterior. Se reescribe como resumen de las dos etapas
% (media página) y remite a los capítulos 4 y 5.
\section{Metodología}

Se utilizarán distintas mediciones reales de $O_3$ provenientes de múltiples estaciones de RAMA/SIMAT junto con su temporalidad provenientes de \textit{RAMA/SIMAT}, sobre la cual se tienen establecidos 5 niveles predefinidos. Estos servirán para distinguir entre simple modificación de una medición y daño sobre la aplicación, los efectos pueden verse con la siguiente función donde \textit{f} es la función que asigna una concentración de $O_3$ a uno de los niveles establecidos por la NOM-172 \cite{nom172_semarnat_2024} , Xreal es la medicion original y Xrecibido la medición recuperada después de la modificación y el descifrado.

\[
D =
\begin{cases}
1, & \text{si } f\left(x_{\text{real}}\right) \neq f\left(x_{\text{recibido}}\right),\\
0, & \text{si } f\left(x_{\text{real}}\right) = f\left(x_{\text{recibido}}\right).
\end{cases}
\]


Las mediciones obtenidas son codificadas en \textit{payload} usando \textit{CayenneLPP}, posteriormente el \textit{payload} se cifra conforme al mecanismo empleado por LoRaWAN, se modifican posiciones específicas del texto cifrado, las modificaciones son seleccionadas de manera sistemática, lo que nos permite aislar el potencial efecto de la posición atacada después de descifrar el \textit{payload} y revisar la medición recibida. 

Como primera etapa de detección se usarán diversos modelos de aprendizaje supervisado como linea base de detección, con la finalidad de determinar cómo cambia la detectabilidad conforme cambia la posición de bit atacada, identificar falsos positivos, falsos negativos y, especialmente, ataques dañinos no detectados.

En una segunda etapa se incorporará el historial temporal de la estación objetivo y el contexto proporcionado por otras estaciones para estimar el comportamiento esperado de la medición; las desviaciones entre el valor esperado y el observado se utilizarán como señal de anomalía.

Finalmente, se evaluará el funcionamiento complementario de ambos mecanismos mediante un esquema de análisis en dos etapas. La primera estará orientada a identificar alteraciones distinguibles mediante aprendizaje supervisado, mientras que la segunda etapa incorporará información temporal y espacial para complementar la detección en aquellos escenarios en los que la información aislada resulte insuficiente.

% REVISAR: falta la contribución principal: NERAD y su evaluación con protocolo
% temporal y registro previo.
\section{Contribución}

\begin{itemize}

    \item Caracterización del daño producido por ataques de \textit{bit flipping} mediante la modificación sistemática de posiciones específicas de bit, relacionando la posición atacada con la magnitud de la alteración, el daño producido y la detectabilidad del ataque.
    
    \item Metodología experimental reproducible para la generación controlada de muestras atacadas, conservando la información correspondiente al valor real, valor recibido, bit atacado, dirección del efecto, daño y severidad.
    
    \item Evaluación del desempeño de mecanismos de detección basados en \textit{ML}, considerando no sólo métricas generales de clasificación, sino también los ataques dañinos que permanecen sin detectar y las mediciones legítimas señaladas incorrectamente como ataques.

\end{itemize}

% REVISAR: ahora son 7 capítulos (4: bit flipping; 5: NERAD, método; 6: resultados;
% 7: conclusiones). Actualizar al final.
\section{Descripción del contenido}

A continuación, se presenta la estructura del presente trabajo de tesis:

\begin{itemize}

    \item[$\bullet$] En el capítulo 2 se establece el sustento teórico necesario para comprender el trabajo, incluyendo conceptos relacionados con \textit{IoT}, \textit{LoRa}, \textit{LoRaWAN}, su arquitectura y mecanismos de seguridad, así como cifrado, maleabilidad, ataques de \textit{bit flipping} y conceptos de aprendizaje automático.

    \item[$\bullet$] En el capítulo 3 se realiza una revisión de trabajos relacionados con vulnerabilidades y ataques de \textit{bit flipping} en \textit{LoRaWAN/LPWAN}, así como de mecanismos de detección basados en \textit{ML}.

    \item[$\bullet$] En el capítulo 4 se presenta la metodología experimental para la detección de ataques sobre datos provenientes de RAMA/SIMAT. Se describen el tratamiento de los datos, la codificación mediante CayenneLPP, el cifrado, el modelo de amenaza y generación de ataques, la definición formal de daño y severidad, la partición temporal, los modelos supervisados, la construcción de series temporales y multiestación, así como los criterios de evaluación.

    \item[$\bullet$] En el capítulo 5 se presentan y analizan los resultados correspondientes a la caracterización daño--detectabilidad, la línea base supervisada, la evaluación temporal-espacial y la comparación del funcionamiento individual y complementario de ambos mecanismos.

    \item[$\bullet$] En el capítulo 6 se responden la hipótesis y los objetivos planteados, se sintetizan las principales contribuciones, se presentan las limitaciones del trabajo y se señalan posibles líneas de trabajo futuro.

\end{itemize}

% =====================================================================
% Capítulo 2. Sustento teórico                     Presupuesto: 14 páginas
% ---------------------------------------------------------------------
% Regla de oro: solo entra la teoría que se USA después. Para cada sección
% pregúntate «¿en qué capítulo la necesito?». Si no hay respuesta, sobra.
% Cada concepto: definición precisa + cita + (si aplica) ecuación o figura.
% No es un libro de texto: no expliques lo que un lector de maestría ya sabe.
% =====================================================================
\chapter{\raggedright \parbox{\textwidth}{Sustento teórico}}
\label{cap:sustento-teorico}

% Breve introducción del capítulo
DESARROLLAR INTRODUCCION AQUI. 

\section{Internet de las Cosas y redes LPWAN}
\label{sec:iot-lpwan}
El Internet de las Cosas (IoT, por sus siglas en inglés, \textit{Internet of Things}) permite la integración de dispositivos como sensores y actuadores~\cite{atzori2010iot}. Los dispositivos se comunican mediante tecnologías de comunicación como Wi-Fi, Bluetooth o las utilizadas en redes celulares. Los sensores recolectan información de variables como temperatura, humedad, concentración de gases u otros parámetros ambientales, la cual puede ser utilizada para monitoreo, automatización y toma de decisiones.

La tecnología de comunicación seleccionada puede no satisfacer los requisitos del despliegue, en términos de costo de infraestructura, consumo energético y autonomía.

Las redes de área amplia y bajo consumo (\textit{LPWAN}, por sus siglas en inglés) están diseñadas para proporcionar comunicación inalámbrica de largo alcance y bajo consumo energético. Estas priorizan la eficiencia energética y la cobertura sobre la tasa de transmisión, por lo que resultan adecuadas para dispositivos que envían pequeños volúmenes de datos de manera periódica. Estas características las hacen relevantes en redes de sensores, donde la información recolectada puede influir en procesos de monitoreo, control o toma de decisiones.

\subsection{Redes de sensores}
\label{subsec:redes-sensores}

La distribución de sensores en distintas ubicaciones permite medir una o varias variables y transmitir las mediciones obtenidas para su procesamiento. En una red de monitoreo atmosférico, cada punto de medición se denomina estación y puede albergar uno o varios sensores ~\cite{akyildiz2002wsn}.

La información generada por un sensor depende de cómo esté implementado; por lo general incluye el valor de la variable, la unidad de medida, el instante de medición y la ubicación o identificación del sensor. Una vez que se dispone de esa información, se almacena y se organiza.

Las mediciones de una misma variable registradas por un sensor a lo largo del tiempo forman una serie temporal. Su organización según el instante de medición permite analizar cómo cambia la variable durante el periodo observado. Cuando se dispone de registros de varios sensores, también es posible comparar las mediciones obtenidas en distintas ubicaciones.

\section{Mediciones y series temporales}
\label{subsec:medicioines-serie-temporal}

Una medición es el resultado de observar una variable de un fenómeno mediante un dispositivo o sistema de medición, en un momento y una ubicación determinados. Los valores registrados pueden variar debido a cambios en el fenómeno observado o a errores o fallas del sistema de medición ~\cite{teh2020sensor}. Por ello, un cambio en las mediciones no implica, por sí solo, un problema en los datos.

El análisis de las mediciones en su orden temporal permite caracterizar la evolución del fenómeno observado. Una serie temporal puede presentar tendencias, patrones periódicos y fluctuaciones asociadas con las condiciones del entorno ~\cite{hyndman2021fpp3}. Estas variaciones pueden formar parte del comportamiento habitual del fenómeno y, por tanto, no constituyen necesariamente anomalías.

Las mediciones de distintas estaciones pueden tener relaciones asociadas con el fenómeno observado, incluso cuando sus valores no sean idénticos. Estas relaciones proporcionan un contexto espacial que complementa la información temporal de cada estación y puede ayudar a identificar desviaciones respecto al comportamiento esperado.

\section{Anomalías en los datos de redes de sensores}
\label{sec:Anomalías-datos-redes-sensores}

Una anomalía en los datos puede describirse como una observación o un conjunto de observaciones que se aparta del comportamiento esperado. Que un valor sea alto o bajo no basta para considerarlo anómalo; es necesario evaluarlo respecto de una referencia que considere el fenómeno observado y su contexto temporal y espacial. Detectar estas desviaciones no demuestra por sí solo que exista una falla o un ataque, por lo que se requiere un análisis adicional para determinar su causa ~\cite{chandola2009anomaly}.

Las anomalías pueden manifestarse tanto en observaciones individuales como en secuencias de mediciones. Algunos valores pueden parecer plausibles cuando se examinan por separado, aunque su comportamiento conjunto se aparte de lo esperado; por ejemplo, cuando presentan una desviación sostenida respecto de una referencia. Por ello, analizar la evolución temporal de las mediciones puede ayudar a identificar alteraciones que pasarían inadvertidas al evaluar cada valor de manera aislada.

% ---------------------------------------------------------------------
% PENDIENTE (orden sugerido; el texto existente de arriba se reacomoda
% cuando lleguemos a este capítulo):
% ---------------------------------------------------------------------
\section{LoRa y LoRaWAN}                                    % 4 págs.
% Debe contestar: arquitectura (ED, GW, NS, AS); claves de sesión de la
% versión 1.1; cómo se cifra el FRMPayload (AES-CTR) y dónde se verifica el
% MIC. Debe tener: figura de la arquitectura (figuras/arquitectura_lorawan.tex)
% marcando el tramo NS -> AS. Cita: especificación 1.1.

\section{Maleabilidad del cifrado y ataques de \textit{bit flipping}}  % 1.5 págs.
% Debe contestar: por qué XOR con un flujo de clave es maleable (ecuación
% C' = C xor Delta  =>  P' = P xor Delta) y por qué el MIC no lo evita en el
% tramo NS -> AS.

\section{Codificación CayenneLPP}                           % 0.5 págs.
% Debe contestar: cómo se codifica una medición (tipo, resolución, bytes) y
% por qué eso fija el efecto de cada posición de bit.

\section{Ozono, la RAMA y el índice de calidad del aire}    % 1.5 págs.
% Debe contestar: por qué el O3 varía tanto (ciclo diario, estacionalidad,
% diferencias entre zonas); qué es la RAMA/SIMAT; bandas de la NOM-172 para
% O3 (58/90/135/175 ppb), que definen el daño en el capítulo 4.

\section{Redes LSTM y ensambles}                            % 2 págs.
% Debe contestar: qué es una LSTM (ecuaciones de las compuertas, breve), por
% qué sirve para series temporales, qué es un ensamble y por qué reduce la
% varianza entre semillas.

\section{Evaluación de detectores de anomalías}             % 2 págs.
% Debe contestar: matriz de confusión; precisión, recall y F1; umbral por
% percentil de datos limpios y su relación con la tasa de falsas alarmas;
% por qué la precisión cae cuando la prevalencia es baja (0.77 %).
% Citas: saito2015pr, chicco2020mcc.

% =====================================================================
% Capítulo 3. Trabajos relacionados                Presupuesto: 6 páginas
% ---------------------------------------------------------------------
% No es una lista de resúmenes. Cada trabajo se cuenta con 3 preguntas:
%   ¿qué hicieron? ¿qué obtuvieron (con número)? ¿en qué difiere de esta tesis?
% Agrupa por tema, no por autor. Termina con una tabla comparativa y un
% párrafo que nombre la BRECHA que cubre la tesis.
% Usa \cite{clave} siempre; nunca cites algo que no hayas leído.
% =====================================================================
\chapter{Trabajos relacionados}
\label{cap:trabajos}

% Introducción del capítulo (3--4 líneas): qué grupos de trabajos se revisan y por qué.

\section{Seguridad de LoRaWAN y ataques de \textit{bit flipping}}
\label{sec:rel-lorawan}
% alizadeh2023bitflipping (ICM): qué datos, cómo generan los ataques (bits
% aleatorios), qué reportan. Diferencia: no relacionan la posición del bit con
% el daño en la decisión. + otros análisis de seguridad de LoRaWAN 1.1.

\section{Detección de anomalías en redes de monitoreo atmosférico}
\label{sec:rel-calidad-aire}
% allka2024paraad (PARAAD) y allka2024stcraad (STC-RAAD): mismo grupo, O3 en
% Barcelona. Qué anomalías sintéticas usan (sesgo, deriva), qué métricas.
% PARAAD es el comparador de esta tesis: decirlo aquí.

\section{Control de calidad con contexto espacial}
\label{sec:rel-espacial}
% lussana2010sct (SCT), baserud2020titan (TITAN), dereszynski2011spatiotemporal,
% deng2021gdn (GDN): cómo usan a los vecinos y qué pasa cuando un vecino
% también está mal (motiva el filtro de NERAD).

\section{Detección de intrusiones en IoT}
\label{sec:rel-ids}
% garciasaez2025fedwla (FedWLA): ataca el tráfico de red, no el dato; la
% prevalencia de ataques (28--49 %) no es comparable con 0.77 %.

\section{Síntesis y brecha}
\label{sec:rel-brecha}
% Tabla \ref{tab:comparacion-trabajos}: filas = trabajos; columnas = datos,
% supervisado, usa vecinas, amenaza, resultado reportado.
% Párrafo final: qué combinación no existe y esta tesis aporta.

% =====================================================================
% Capítulo 4. Etapa 1: bit flipping en LoRaWAN 1.1  Presupuesto: 9 páginas
% ---------------------------------------------------------------------
% Este capítulo tiene método Y resultados de la etapa 1 (es una etapa cerrada).
% Su función en la historia: demostrar que (1) un atacante sin la clave puede
% causar daño real con valores plausibles, (2) la medición aislada no basta y
% (3) el contexto temporal-espacial ayuda. Eso motiva la etapa 2.
% Fuentes en el repo: docs/01_umbrales_normativos.md, docs/02_codificacion.md,
% docs/protocolo_lstm_v4_multianual.md, results/lstm_v4_multianual_v1/,
% experiments/cayenne_mod/cierre_distribucion_pesos.md.
% =====================================================================
\chapter{Etapa 1: daño y detectabilidad del \textit{bit flipping}}
\label{cap:bit-flipping}

\section{Modelo de amenaza}
\label{sec:amenaza}
% Debe contestar: quién es el atacante, dónde está (tramo NS -> AS), qué sabe
% (estructura CayenneLPP), qué NO sabe (claves), qué quiere (cambiar la banda
% comunicada). Figura: figuras/modelo_amenaza.tex.

\section{Generación sistemática de ataques}
\label{sec:generacion}
% Debe contestar: cómo se codifica, cifra, modifica y descifra; qué posiciones
% de bit se recorren; qué se guarda de cada muestra (valor real, recibido, bit,
% dirección, daño). Debe permitir que otro lo reproduzca.

\section{Definición del daño}
\label{sec:dano}
% Ecuación D = 1[f(x_real) != f(x_recibido)] con f = banda NOM-172.
% Ejemplo: 32 ppb lejos del corte no dañan; 2 ppb junto al corte sí.

\section{Detección: medición aislada contra contexto}
\label{sec:deteccion-e1}
% Capa 1 (medición aislada) y LSTM con contexto espacial; detector V4
% (31 LSTM, ajuste 2020--2022, calibración 2023, prueba 2024--2026).

\section{Resultados}
\label{sec:resultados-e1}
% - Tabla: detección con 32 y 64 ppb, aislada contra LSTM (estación CCA, 2025).
% - Óptimo interior del atacante: máximo daño no detectado en medio.
% - Bits 3--5 concentran el daño no detectado.
% - Deriva estacional -> perfil causal de 14 días (se hereda a la etapa 2).
% - Contramedida en CayenneLPP: resultado negativo (255 -> 567 cambios de
%   banda no detectados en 179 399 intentos).

\section{De la etapa 1 a la etapa 2}
\label{sec:puente}
% Media página: qué se aprendió y por qué el problema se generaliza a
% «qué estación reporta datos anómalos», sin importar la causa.

% =====================================================================
% Capítulo 5. NERAD: datos, diseño y protocolo     Presupuesto: 14 páginas
% ---------------------------------------------------------------------
% La metodología debe permitir que otro investigador REPITA el trabajo.
% Para cada decisión: qué se hizo, con qué parámetros y POR QUÉ.
% Debe contestar:
%  - ¿Con qué datos? (fuente, años, estaciones, limpieza, faltantes)
%  - ¿Qué anomalías y cómo se inyectan? (4 escenarios, calendarios pareados)
%  - ¿Cómo se evita la fuga de información? (división temporal, registro previo)
%  - ¿Cómo funciona el detector? (ecuaciones, no solo prosa)
%  - ¿Contra qué se compara y por qué es justo? («peras con peras»)
%  - ¿Cómo se mide el éxito? (métricas, umbral, intervalos)
% Escribir en pasado («se entrenó»), voz impersonal.
% =====================================================================
\chapter{NERAD: datos, diseño y protocolo experimental}
\label{cap:metodo}

\section{Datos de la RAMA}
\label{sec:datos}
% Estaciones, años usados, O3 horario, criterio de >= 12 h con dato por día,
% vecinas más cercanas (4), datos validados de SIMAT para ago--sep 2026.

\section{Escenarios de anomalía}
\label{sec:escenarios}
% Tabla: 1. proporcional (+-delta 20/50/100 %), 2. escalón aditivo (+-10/20/40 ppb),
% 3. ataque coordinado a 0/1/2/4 vecinas, 4. sesgo sostenido (+20/+40 ppb, 8 semanas).
% Episodios de 1 a 7 días desde las 00 h; sentido sorteado; calendarios pareados.

\section{Protocolo temporal y registro previo}
\label{sec:protocolo}
% Figura: figuras/protocolo.tex (línea de tiempo del coloquio, adaptada).
% 2020--2021 entrenamiento; 2022 selección; 2023 calibración sin ataques;
% 2024--jul 2026 banco (30 calendarios); ago--sep 2026 despliegue (descriptivo).
% Registro previo, doble corrida, huellas SHA-256.

\section{Diseño incremental del detector}
\label{sec:escalera}
% La escalera L0 único -> Ens-V0-R -> Ens-Vb-R -> NERAD: qué agrega cada
% escalón y qué pregunta respondió. Fuente: docs/redaccion_escalera_ablacion.md.

\section{Arquitectura de NERAD}
\label{sec:arquitectura}
% Ecuaciones: z diario; regla de vecina sospechosa (mediana de las otras 3,
% tau = 0.2596, mismo signo; con 3 o 4 sospechosas no se excluye ninguna);
% base vecinal; ventana 24 h x 17 columnas; ensamble de 5 LSTM; decisión diaria.
% Figura: figuras/arquitectura.tex.

\section{Comparador: PARAAD-SA}
\label{sec:comparador}
% Reimplementación adaptada a la RAMA, entrenada sin etiquetas con 2020--2022
% y calibrada igual. Declarar que no es el código original.

\section{Métricas y criterios de comparación}
\label{sec:metricas}
% Unidad: estación-día. Umbral p99 de 2023 limpio. Recall, precisión, F1,
% días de falsa alarma al año por estación. Bootstrap por bloques estación x mes.

% =====================================================================
% Capítulo 6. Resultados y discusión               Presupuesto: 14 páginas
% ---------------------------------------------------------------------
% Cada sección sigue el mismo patrón (párrafo por paso):
%   1. Qué pregunta responde el experimento (una oración).
%   2. La figura o tabla, citada en el texto ANTES de que aparezca.
%   3. Qué se ve, con los números clave (no repetir toda la tabla).
%   4. Qué significa y qué NO se puede afirmar.
% Los resultados negativos también se reportan.
% Cada número debe salir de un CSV del repo; anota la fuente en el comentario
% de la figura.
% =====================================================================
\chapter{Resultados y discusión}
\label{cap:resultados}

\section{Elección del diseño (2022)}
\label{sec:res-escalera}

\section{Elección del punto de operación}
\label{sec:res-punto}

\section{Una estación atacada: NERAD contra PARAAD-SA}
\label{sec:res-una}
% Figura de prueba de la paleta (borrar o conservar cuando lleguemos aquí):
\begin{figure}[htbp]
  \centering
  % figuras/recall_banco_p99.tex
% Recall con una estación atacada (escenario E1), umbral p99 de 2023.
% Fuente: results/linea_base/rev24/metricas_banco.csv, filas E1, nivel 0.01,
%   detectores Ens-Vb-X0.2596 (NERAD) y PARAAD. Verificado el 2026-11-27.
\begin{tikzpicture}
\begin{axis}[
  ejes tesis,
  width=0.9\textwidth, height=6cm,
  ybar, bar width=9pt,
  ymin=0, ymax=1.05, ytick={0,0.2,...,1},
  ylabel={Recall},
  xlabel={Magnitud de la anomalía},
  symbolic x coords={d20,d50,d100,a10,a20,a40},
  xtick=data,
  xticklabels={$\pm$20\,\%, $\pm$50\,\%, $\pm$100\,\%, $\pm$10 ppb, $\pm$20 ppb, $\pm$40 ppb},
  enlarge x limits=0.1,
  legend style={at={(0.02,0.98)}, anchor=north west},
  legend cell align=left,
]
\addplot[barra=ensX] coordinates {(d20,0.2020) (d50,0.8019) (d100,0.9319) (a10,0.2279) (a20,0.8312) (a40,0.9584)};
\addlegendentry{NERAD}
\addplot[barra=paraad] coordinates {(d20,0.0884) (d50,0.3647) (d100,0.6753) (a10,0.1013) (a20,0.5656) (a40,0.9201)};
\addlegendentry{PARAAD-SA}
\end{axis}
\end{tikzpicture}

  \caption[Recall en el banco de pruebas con umbral p99]{Recall de NERAD y
  PARAAD-SA con una estación atacada en el banco de pruebas (2024 a julio de
  2026, 30 calendarios, 1\,777 estación-días atacados), con el umbral p99
  calibrado en 2023 sin ataques.}
  \label{fig:recall-banco-p99}
\end{figure}

\section{Despliegue en agosto y septiembre de 2026}
\label{sec:res-despliegue}

\section{Ataque coordinado a las vecinas}
\label{sec:res-coordinado}

\section{Sesgo sostenido}
\label{sec:res-sesgo}

\section{Resultados negativos}
\label{sec:res-negativos}
% Examen 2017--2019, EP1, T0, fusión con PARAAD-SA.

\section{Exploración con imágenes TINTO}
\label{sec:res-tinto}

\section{Amenazas a la validez}
\label{sec:validez}
% Anomalías sintéticas; el banco orientó revisiones; despliegue descriptivo;
% reimplementación de PARAAD-SA; ventana de despliegue fuera de la temporada alta.

% =====================================================================
% Capítulo 7. Conclusiones                         Presupuesto: 4 páginas
% ---------------------------------------------------------------------
% Aquí no hay resultados nuevos ni citas nuevas. Debe contestar:
%  1. ¿Se cumple la hipótesis? Sí / no / en parte, con el número que lo dice.
%  2. Cada objetivo específico, uno por uno, y dónde se cumplió (\ref).
%  3. Contribuciones (las mismas del capítulo 1, ahora demostradas).
%  4. Limitaciones (honestas, sin repetir todo el capítulo 6).
%  5. Trabajo futuro concreto (no «mejorar el modelo»).
% =====================================================================
\chapter{Conclusiones}
\label{cap:conclusiones}

\section{Respuesta a la hipótesis}
\section{Cumplimiento de los objetivos}
\section{Contribuciones}
\section{Limitaciones}
\section{Trabajo futuro}


% =====================
% Referencias
% =====================
\printbibliography

\end{document}
